\chapter{Experimentos y Resultados}
\label{cap:experimentos_resultados}

\section{Experimentos}
\label{sec:experimentos}
En este capítulo se evalúa el comportamiento fuera de muestra de las estrategias construidas con la metodología del Capítulo~\ref{cap:metodologia}. El objetivo no es solo comparar rentabilidad ajustada por riesgo, sino también analizar el coste computacional de cada enfoque para estudiar el compromiso entre calidad de solución y tiempo de ejecución.

\subsection{Marco Experimental}

La evaluación se realiza mediante un esquema \textit{walk-forward} con ventanas deslizantes de 756 días (aprox. 3 años) para entrenamiento y 21 días (aprox. 1 mes) para test, evaluando 18 ventanas en total. En cada ventana, los pesos se estiman con datos de entrenamiento y se evalúan exclusivamente en el bloque temporal posterior, evitando fuga de información y reproduciendo un uso operativo realista.

Se comparan cinco estrategias: \textbf{HRP Cuántico (QHRP)} (ordenación jerárquica a partir de distancias entre matrices de densidad); \textbf{HRP Clásico} (ordenación jerárquica basada en distancia de correlación); \textbf{HRP Kernel} (ordenación jerárquica basada en distancia MMD gaussiana); \textbf{Markowitz} (asignación proporcional al inverso de la varianza diagonal); y \textbf{Pesos Iguales (1/N)} (cartera equiponderada como línea base).

Para cada método y cada ventana se computan: \textbf{Sharpe ratio anualizado} (medida principal de rendimiento ajustado por riesgo), \textbf{PSR (Probabilistic Sharpe Ratio)} (probabilidad de superar un Sharpe objetivo), y \textbf{MinTRL} (indicador de robustez estadística de la señal). Adicionalmente, se agregan todos los retornos fuera de muestra por método para obtener un Sharpe global acumulado.


\section{Resultados}
\label{sec:resultados}

\begin{table}[Rendimiento walk-forward]{TB:WALKFORWARD}{Rendimiento walk-forward por ventana (18 ventanas).}
\begin{tabular}{lccc}
\hline
\textbf{Método} & \textbf{Sharpe medio} & \textbf{Sharpe desv.} & \textbf{PSR medio} \\
\hline
HRP Cuántico & 1.5016 & 3.7837 & 0.6191 \\
HRP Clásico & 1.2754 & 3.9986 & 0.5976 \\
HRP Kernel & 1.2928 & 3.7140 & 0.6054 \\
Markowitz & 1.3449 & 3.7788 & 0.6089 \\
Pesos Iguales & 1.5076 & 3.9906 & 0.6113 \\
\hline
\end{tabular}
\end{table}

La Tabla~\ref{TB:WALKFORWARD} muestra que la diferencia entre QHRP (1.5016) y Pesos Iguales (1.5076) es pequeña frente a su variabilidad interventana (desviaciones cercanas a 4), por lo que en una ventana concreta cualquiera puede quedar por encima. Esto significa que el comportamiento mensual está fuertemente influido por cambios de régimen y ruido de estimación, lo que sugiere que las diferencias no son estadísticamente robustas a nivel mensual y deben interpretarse con cautela. Por tanto, esta tabla sirve para evaluar robustez local, pero no basta por sí sola para decidir el mejor método global.

\begin{table}[Sharpe agregado OOS]{TB:SHARPEAGREGADO}{Sharpe agregado fuera de muestra (18 ventanas).}
\begin{tabular}{lc}
\hline
\textbf{Método} & \textbf{Sharpe OOS agregado} \\
\hline
HRP Cuántico & 1.4441 \\
HRP Clásico & 1.2378 \\
HRP Kernel & 1.1803 \\
Markowitz & 1.2518 \\
Pesos Iguales & 1.3193 \\
\hline
\end{tabular}
\end{table}

La Tabla~\ref{TB:WALKFORWARD} y la Tabla~\ref{TB:SHARPEAGREGADO} muestran que QHRP lidera en términos agregados, mientras que la variabilidad mensual impide declarar un ganador estadísticamente robusto a corto plazo. En consecuencia, la comparación debe hacerse en dos niveles: estabilidad local por ventana y rendimiento global acumulado. La Tabla~\ref{TB:COSTECOMPUTACIONAL} completa la lectura al mostrar que la mejora financiera de QHRP viene acompañada de un coste computacional muy superior.

\begin{table}[Coste computacional]{TB:COSTECOMPUTACIONAL}{Coste computacional medio por ventana ($N=169$, $P=6$).}
\begin{tabular}{lcc}
\hline
\textbf{Etapa / Método} & \textbf{Media (s)} & \textbf{Desv. (s)} \\
\hline
Matrices densidad (cuántico) & 1903.279 & 290.486 \\
Distancia Frobenius (cuántico) & 6.284 & 19.258 \\
Ordenación cuántica (Ward) & 0.004 & 0.005 \\
Ordenación clásica (correlación) & 0.069 & 0.015 \\
Ordenación kernel (MMD) & 510.007 & 119.269 \\
Covarianza y pesos & 1.536 & 0.290 \\
\hline
TOTAL HRP Cuántico & 1909.566 & 291.434 \\
TOTAL HRP Clásico & 0.069 & 0.015 \\
TOTAL HRP Kernel & 510.007 & 119.269 \\
\hline
\end{tabular}
\end{table}

La Tabla~\ref{TB:COSTECOMPUTACIONAL} muestra que el cuello de botella de QHRP está en la simulación de las matrices de densidad, con un coste aproximadamente $2.8\times 10^4$ veces superior al clásico. La fase de ordenación es marginal frente a esa simulación, así que la limitación real del método es su escalabilidad en entornos de rebalanceo frecuente.

\subsection{Relación entre estructura visual y desempeño}
La comparación de matrices ordenadas y no ordenadas del Capítulo~\ref{cap:metodologia} ayuda a interpretar por qué en algunos casos el panel derecho ofrece mejor desempeño que el izquierdo. El panel ordenado concentra activos similares en bloques cuasi-diagonales, lo que favorece particiones más coherentes en la bisección recursiva de HRP. En consecuencia, se reduce la mezcla de activos poco afines dentro de subcarteras y mejora la asignación de riesgo.

No obstante, esta relación no es determinista: una estructura visual más definida no garantiza necesariamente un mejor rendimiento financiero. Sin embargo, sí indica que el método de distancia está capturando estructura jerárquica relevante y económicamente significativa.

\subsection{Robustez frente a permutaciones de activos}

Para validar que la estructura jerárquica es invariante al orden inicial de los activos, se realizó un experimento permutando aleatoriamente el orden de entrada en cinco ocasiones. La Tabla~\ref{TB:FIVEPERMS} resume los resultados:

\begin{table}[Cinco permutaciones de activos]{TB:FIVEPERMS}{Resultados del test con cinco desordenaciones distintas de activos.}
\begin{tabular}{cccccc}
\hline
\textbf{perm} & \textbf{seed} & \textbf{gain\_classical} & \textbf{gain\_quantum} & \textbf{sim\_classical} & \textbf{sim\_quantum} \\
\hline
1 & 11 & 0.1402 & 0.0234 & 1.0000 & 1.0000 \\
2 & 22 & 0.1479 & 0.0314 & 1.0000 & 1.0000 \\
3 & 33 & 0.1421 & 0.0253 & 1.0000 & 1.0000 \\
4 & 44 & 0.1448 & 0.0280 & 1.0000 & 1.0000 \\
5 & 55 & 0.1472 & 0.0307 & 1.0000 & 1.0000 \\
\hline
\textbf{Media} & -- & \textbf{0.1444} & \textbf{0.0278} & \textbf{1.0000} & \textbf{1.0000} \\
\hline
\end{tabular}
\end{table}

La similitud 1.0000 en todas las corridas indica que la estructura de bloques se recupera idénticamente independientemente del orden inicial de activos. Los \textit{gains} positivos confirman que ambos métodos mejoran consistentemente. La Figura~\ref{FIG:PERMUTACIONES_CORR} y Figura~\ref{FIG:PERMUTACIONES_QUANTUM} muestran que en matrices sin ordenar la estructura se pierde, pero tras aplicar HRP/QHRP se recuperan patrones cuasi-diagonales equivalentes. Esto valida que los resultados no son artefactos visuales sino consecuencias del procedimiento jerárquico de ordenación.

\begin{figure}[Invarianza frente a permutaciones: correlación]{FIG:PERMUTACIONES_CORR}{Ejemplo: matrices de correlación sin ordenar y tras HRP clásico en una permutación representativa de activos.}
	\image{9cm}{}{fig3-shuffled}
\end{figure}

\begin{figure}[Invarianza frente a permutaciones: cuántica]{FIG:PERMUTACIONES_QUANTUM}{Ejemplo: matrices de distancia cuántica sin ordenar y tras QHRP en una permutación representativa de activos.}
	\image{9cm}{}{fig4-shuffled}
\end{figure}

\subsection{Comparación con el paper original}
Para evaluar el grado de reproducibilidad, se comparan los resultados de este trabajo con los reportados por el paper original de QHRP en tres dimensiones: (i) ranking relativo entre métodos, (ii) magnitud de Sharpe fuera de muestra y (iii) comportamiento visual de las matrices ordenadas.

En este trabajo, el ranking por Sharpe agregado fuera de muestra (Tabla~\ref{TB:SHARPEAGREGADO}) es:

\begin{center}
QHRP (1.4441) $>$ Pesos Iguales (1.3193) $>$ Markowitz (1.2518) $>$ HRP Clásico (1.2378) $>$ HRP Kernel (1.1803).
\end{center}

En términos de reproducibilidad, el criterio principal no es la coincidencia exacta en valores absolutos, sino la consistencia cualitativa: que QHRP mantenga un desempeño competitivo frente al HRP clásico y que la ordenación cuántica siga revelando estructura jerárquica coherente en las matrices reordenadas.

\begin{table}[Comparación con paper original (Tabla 1)]{TB:COMPARACIONPAPER_T1}{Comparación de métricas por ventana con la Tabla 1 del paper original.}
\begin{tabular}{lccc}
\hline
\textbf{Método} & \textbf{Sharpe (est./pap.)} & \textbf{Desv. (est./pap.)} & \textbf{PSR (est./pap.)} \\
\hline
QHRP & 1.5016 / 1.4658 & 3.7837 / 3.7234 & 0.6191 / 0.6172 \\
HRP Clásico & 1.2754 / 1.2190 & 3.9986 / 3.8598 & 0.5976 / 0.5970 \\
HRP Kernel & 1.2928 / 1.2944 & 3.7140 / 3.6671 & 0.6054 / 0.6055 \\
Markowitz & 1.3449 / 1.3474 & 3.7788 / 3.7785 & 0.6089 / 0.6093 \\
Pesos Iguales & 1.5076 / 1.3565 & 3.9906 / 3.9065 & 0.6113 / 0.6043 \\
\hline
\end{tabular}
\end{table}

La Tabla~\ref{TB:COMPARACIONPAPER_T1} revela una cercanía entre Sharpe, desviación y PSR respecto al paper que sugiere que el pipeline experimental está bien reproducido a nivel de ventana. Las discrepancias más visibles en algunos métodos (por ejemplo, Pesos Iguales) son compatibles con sensibilidad al periodo de mercado más que con un fallo metodológico. En consecuencia, la réplica puede considerarse fiel en dinámica estadística general.

\begin{table}[Comparación con paper original (Tabla 2)]{TB:COMPARACIONPAPER_T2}{Comparación de Sharpe agregado fuera de muestra con la Tabla 2 del paper original.}
\begin{tabular}{lcc}
\hline
\textbf{Método} & \textbf{Este trabajo} & \textbf{Paper original} \\
\hline
QHRP & 1.4441 & 1.3975 \\
HRP Clásico & 1.2378 & 1.1378 \\
HRP Kernel & 1.1803 & 1.2186 \\
Markowitz & 1.2518 & 1.2557 \\
Pesos Iguales & 1.3193 & 1.1983 \\
\hline
\end{tabular}
\end{table}

Como muestran las Tablas~\ref{TB:COMPARACIONPAPER_T1} y~\ref{TB:COMPARACIONPAPER_T2}, el pipeline es reproducible en dinámica estadística y ranking global, con variaciones menores en el orden intermedio debido a la sensibilidad de la muestra.

\begin{figure}[Comparación visual Fig. 3 (este trabajo)]{FIG:COMP_PAPER_FIG3_THIS}{Figura de correlaciones reordenadas obtenida en este trabajo con el mismo esquema metodológico.}
	\image{12.5cm}{}{fig3-corr-matrices}
\end{figure}

La Figura~\ref{FIG:COMP_PAPER_FIG3_THIS} muestra que al recuperar bloques y separación entre grupos, se valida que la implementación reproduce la lógica económica del paper (activos similares quedan próximos en el ordenamiento). Las diferencias finas de contraste o tamaño son esperables por cambios en muestra y normalización, y no invalidan la réplica metodológica.

\begin{figure}[Comparación visual Fig. 4 (este trabajo)]{FIG:COMP_PAPER_FIG4_THIS}{Figura de distancias cuánticas obtenida en este trabajo.}
	\image{12.5cm}{}{fig4-best-match}
\end{figure}

La Figura~\ref{FIG:COMP_PAPER_FIG4_THIS} confirma que la presencia de bloques también en esta réplica muestra que QHRP mantiene poder de organización estructural fuera del contexto exacto del paper. La no coincidencia celda a celda es secundaria; lo relevante es conservar la separación jerárquica entre grupos de activos.

La comparación visual se conserva como validación cualitativa, pero basta con las figuras reproducidas en este trabajo para mostrar el patrón cuasi-diagonal y la separación jerárquica; no hace falta duplicar las imágenes del paper original. Para esta parte visual se adopta un criterio estructural:

\begin{itemize}
	\item \textbf{Compatibilidad jerárquica:} agrupaciones estables y coherentes en la ordenación.
	\item \textbf{Separación visual:} bloques de alta similitud interna y menor similitud entre bloques.
	\item \textbf{Consistencia:} alineación entre estructura visual y resultados cuantitativos (Sharpe/PSR).
\end{itemize}

Bajo estos criterios, las imágenes obtenidas en este trabajo son comparables a las del paper en términos metodológicos, aunque presenten diferencias estéticas o de granularidad local. Las discrepancias pueden explicarse principalmente por: (i) diferencias en universo temporal (tickers, fechas, regímenes de mercado) que alteran correlaciones y volatilidades, (ii) decisiones de preprocesado y parametrización del modelo ($\alpha$, normalización, criterio MMD), y (iii) detalles de implementación jerárquica (método de enlace, tratamiento de empates en dendrogramas).

Por tanto, si los valores absolutos difieren pero se conserva el patrón cualitativo, la reproducción puede considerarse satisfactoria desde el punto de vista metodológico.

\subsection{Piloto en hardware IBM con 40 activos}
Para reforzar la validación experimental en entorno real, se diseñó un piloto en hardware IBM con una pregunta operativa explícita: dado un subconjunto fijo de activos y el mismo pipeline QHRP, \textbf{¿se conserva la mejora estructural que aporta la ordenación cuántica al pasar de simulación ideal a hardware real?}

La conservación se mide con la métrica \textit{block contrast} sobre la matriz de distancias \(D\). Para \(b=\max(3,\lfloor N/10\rfloor)\), se definen \(\mathcal{N}_b=\{(i,j):0<|i-j|\le b\}\) y \(\mathcal{F}_b=\{(i,j):|i-j|>b\}\); la condición \(0<|i-j|\) excluye explícitamente la diagonal. Sean \(\overline{D}_{\mathcal{N}_b}=\operatorname{mean}\{D_{ij}:(i,j)\in\mathcal{N}_b\}\) y \(\overline{D}_{\mathcal{F}_b}=\operatorname{mean}\{D_{ij}:(i,j)\in\mathcal{F}_b\}\). Entonces, \(\mathrm{BC}(D;b)=\overline{D}_{\mathcal{F}_b}-\overline{D}_{\mathcal{N}_b}\), con \(\Delta\mathrm{BC}=\mathrm{BC}(D_{\text{ord}};b)-\mathrm{BC}(D_{\text{raw}};b)\). Un valor positivo de \(\Delta\mathrm{BC}\) indica que la ordenación incrementa la separación entre pares cercanos y lejanos al eje diagonal.

El piloto se acotó a \(N=40\) activos y \(T=90\) días recientes. Esta decisión responde a restricciones prácticas de hardware (cola, coste de medición y ruido), por lo que se priorizó un tamaño manejable pero representativo. En consistencia con la definición del tensor en la Sección~\ref{sec:construccion_tensor}, \(P=6\) denota el número de características codificadas en el \textit{feature map} (un qubit por característica).

La selección de activos se hizo de forma reproducible y en dos etapas:
\begin{enumerate}
\item \textbf{Criterio principal (liquidez):} ranking descendente por mediana de volumen monetario diario \((\text{precio}\times\text{volumen})\) en la ventana de 90 días. Se usa la mediana, y no la media, por su mayor robustez frente a picos transitorios y valores atípicos.
\item \textbf{Criterio de respaldo (volatilidad):} si tras el filtro de datos faltan activos para completar 40, se añaden los no seleccionados con mayor desviación estándar de retornos en la misma ventana temporal.
\end{enumerate}

La simulación actúa como \textit{baseline}. El pipeline se mantiene idéntico en preprocesado, subconjunto de activos, \textit{feature map}, ordenación jerárquica y criterio de evaluación. Solo cambia la capa de ejecución cuántica (observables ideales en simulador frente a distribuciones medidas en hardware). También se fija \(\alpha=2.0\) en ambos casos para mantener comparabilidad con el flujo principal y conservar una escala angular estable en la construcción de distancias.

\begin{table}[Piloto hardware vs simulación]{TB:PILOTO_BC}{Resultados de \textit{block contrast} en el piloto \(N=40\), \(T=90\), \(P=6\).}
\begin{tabular}{lccc}
\hline
\textbf{Caso} & \textbf{$BC_{\mathrm{raw}}$} & \textbf{$BC_{\mathrm{ord}}$} & \textbf{$\Delta\mathrm{BC}$} \\
\hline
Simulación (baseline) & -0.0037 & 0.0748 & 0.0786 \\
Hardware IBM (256 \textit{shots}/circuito) & -0.0006 & 0.0096 & 0.0102 \\
\hline
\end{tabular}
\end{table}

Los resultados de hardware corresponden a una única ejecución (sin promedio sobre \textit{runs} independientes). La mejora estructural se conserva en signo, pero se atenúa en magnitud: \(\Delta\mathrm{BC}_{\text{hw}}/\Delta\mathrm{BC}_{\text{sim}}\approx 0.1298\), equivalente a una degradación relativa de \(87.0\%\). Esta degradación combina dos efectos: (i) varianza de muestreo por \textit{shots} finitos y (ii) ruido físico del hardware (errores de puerta, lectura y deriva de calibración). Dado que se trata de una sola corrida en hardware, este piloto no permite aislar cuantitativamente ambos efectos ni sostener una atribución causal fina entre ellos.

El criterio visual (bloques cuasi-diagonales en mapas de calor) se utiliza como evidencia cualitativa complementaria, no como prueba suficiente por sí sola. Por ende, el piloto se interpreta como evidencia de \textbf{transferencia metodológica parcial}: la estructura jerárquica de QHRP persiste en hardware real, aunque con menor contraste que en simulación ideal.

Por último, para asegurar reproducibilidad operativa, se guardó una copia local de la matriz de distancias de hardware y se habilitó recuperación desde trabajos ya completados en la plataforma de IBM. Esto permite rehacer figuras comparativas sin relanzar ejecuciones costosas tras reiniciar el kernel.

\subsection{Análisis de segmentación jerárquica: coherencia económica de los grupos}
\label{subsec:analisis_segmentacion}

Más allá del rendimiento financiero, resulta de interés evaluar si los métodos HRP (clásico y cuántico) identifican grupos de activos que tengan sentido económico y sean estables frente a cambios de régimen. Este análisis complementa la evaluación de Sharpe, ya que una segmentación coherente constituye una base sólida para una optimización de carteras robusta.

\subsubsection{Metodología: pureza y entropía sectorial}
\label{subsubsec:metodologia_pureza_entropia}

Para cada grupo detectado por la jerarquía, se calculan dos métricas de composición:

\begin{itemize}
	\item \textbf{Purity (Pureza):} fracción de activos que pertenecen al sector dominante en ese grupo. Un valor cercano a 1.0 indica homogeneidad sectorial; valores bajos (~0.13) indican mezcla de múltiples sectores.
	
	\item \textbf{Entropy (Entropía de Shannon):} medida de diversidad sectorial dentro del grupo. $H = -\sum_i p_i \log p_i$, donde $p_i$ es la proporción de activos en el sector $i$. Valores cercanos a 0 indican uniformidad; valores altos (~0.86) indican que los activos están distribuidos uniformemente entre muchos sectores.
\end{itemize}

Estas métricas permiten evaluar objetivamente si un método agrupa activos que comparten características macroeconómicas o si, por el contrario, mezcla indiscriminadamente factores de riesgo distintos.

\subsubsection{HRP Clásico: Grupo principal más grande}
\label{subsubsec:hrp_clasico_megagrupo}

En el análisis del HRP clásico, el grupo más grande contiene 147 activos (de 169 totales, con 23 sectores distintos, Purity=0.129, Entropy=0.866), revelando un problema fundamental de la ordenación clásica: agrupa un ``mega-grupo'' heterogéneo que contiene simultáneamente:

\begin{itemize}
	\item \textbf{Renta variable:} acciones tecnológicas (AAPL, MSFT, NVDA), industriales (BA, CAT, GE), retail (TJX, HD, NKE) y financieras (JPM, WFC, BAC).
	
	\item \textbf{Renta fija:} 15 ETFs de bonos (AGG, BND, TLT, LQD, SCHZ, VCIT...) que, en condiciones normales, están negativamente correlacionados con la renta variable.
	
	\item \textbf{Commodities:} 9 ETFs de materias primas (DBC, PDBC, GLD, SLV, COPX...) que actúan como coberturas ante inflación.
\end{itemize}

Aunque la métrica de correlación histórica agrupa estos activos bajo un único dendograma, sus dinámicas económicas son heterogéneas. Cuando ocurren shocks de mercado o cambios de régimen, las correlaciones convergen hacia 1.0 (risk-off) o divergen según el factor de riesgo prevalente (ej., rentas variables vs. bonos). En consecuencia, este mega-grupo es frágil para optimización de carteras, ya que su estructura interna puede colapsar rápidamente ante perturbaciones externas.

\subsubsection{HRP Clásico: Grupo crypto}
\label{subsubsec:hrp_clasico_crypto}

En contraste, el grupo de criptomonedas (9 activos: BTC, ETH, SOL, BNB, AVAX, XRP, DOGE, ADA, DOT) presenta estadísticas ideales (Purity=1.0, Entropy=0.0, sector único), indicando que los métodos jerárquicos sí pueden captar agrupamientos significativos cuando existe un factor de riesgo uniforme y dominante. El grupo crypto actúa como validación de que una Purity~=~1.0 es alcanzable y económicamente interpretable.

\subsubsection{QHRP Cuántico: Grupo de Diversificadores/Coberturas}
\label{subsubsec:qhrp_diversificadores}

En el análisis del QHRP, uno de los grupos principales (31 activos con 12 sectores, Purity=0.290, Entropy=0.845) corresponde a una familia de diversificadores y coberturas. A diferencia del mega-grupo del HRP clásico, este grupo está explícitamente compuesto por activos que tienen funciones de cobertura complementarias:

\begin{itemize}
	\item \textbf{Renta fija (9 activos):} AGG, BND, TLT, LQD, SCHZ, VCIT... Actúan como ancla de valor cuando caen las rentas variables (comportamiento anti-cíclico).
	
	\item \textbf{Materias primas (7 activos):} DBC, PDBC, GLD, SLV, COPX, KRBN, PPLT. Generan retornos desacoplados de rentas variables y actúan como coberturas ante inflación.
	
	\item \textbf{Renta variable de mercados emergentes:} VWO, EEM, IEMG. Valuaciones bajas y correlación reducida con rentas variables estadounidenses; tienden a recuperarse primero tras ventas de pánico.
	
	\item \textbf{Outliers estratégicos:} PFE (farmacéutico), BA (defensa), AMD (semiconductores), OXY y SLB (oil E\&P). Cada uno tiene dinámicas propias que los hacen útiles en carteras diversificadas.
\end{itemize}

La clave es que QHRP no agrupa por similitud sectorial estricta, sino por \textbf{coherencia de factores de riesgo}. Estos 31 activos co-varían internamente como bloque de ``protección contra rentas variables'', pero se mueven opuestamente a la renta variable pura en escenarios de estrés.

\subsubsection{QHRP Cuántico: Grupo de Renta Variable Pura}
\label{subsubsec:qhrp_equities}

Otro grupo principal de QHRP (25 activos con Purity=0.400, Entropy=0.825) agrupa la renta variable sistemática. Este grupo contiene:

\begin{itemize}
	\item \textbf{ETFs amplios (14 activos):} SPY, VOO, IVV, QQQ, VTI, EFA, VEA, SPDW, SCHF, ACWI... Con correlaciones ~0.99 entre ellos.
	
	\item \textbf{Acciones individuales (11 activos):} XOM, VLO (Energy), PEP (Consumer), MCD (Retail), HON, HD, BYDDF, XLB, XLK, XLV, XLU, XLY... Componentes de los índices que replican los ETFs amplios.
	
	\item \textbf{Materias primas de riesgo sistémico (2 activos):} USO, UGA (oil, agriculture). En el período analizado, mostraron correlación positiva con rentas variables.
\end{itemize}

El factor común es que todos responden al ``apetito por riesgo global'': cuando sube la tolerancia al riesgo, todos suben; cuando hay pánico, todos caen. La Purity más baja (0.40) es aceptable porque el grupo está definido por su respuesta macroeconómica común, no por homogeneidad sectorial.

\subsubsection{Análisis multinivel: Granularidad creciente}
\label{subsubsec:analisis_multinivel}

El análisis jerárquico permite refinar la segmentación en tres niveles:

\begin{table}[Segmentación multinivel en QHRP]{TB:QHRP_MULTILEVEL}{Estadísticas de segmentación a diferentes niveles de granularidad en QHRP.}
\begin{tabular}{lccccc}
\hline
\textbf{Nivel} & \textbf{Grupos} & \textbf{Tamaño medio} & \textbf{Purity media} & \textbf{Entropy media} & \textbf{Sectores/grupo} \\
\hline
MAIN & 11 & 15.3 & 0.273 & 0.941 & 8.3 \\
SUB & 31 & 5.3 & 0.448 & 0.863 & 3.7 \\
MICRO & 45 & 3.2 & 0.504 & 0.859 & 2.6 \\
\hline
\end{tabular}
\end{table}

La tendencia es clara: a medida que aumenta la granularidad (de MAIN a SUB a MICRO), los grupos se vuelven más homogéneos (Purity sube de 0.27 a 0.50), más compactos (tamaño baja de 15.3 a 3.2) y más especializados (sectores/grupo baja de 8.3 a 2.6). 

Esta progresión indica que la segmentación cuántica es jerárquicamente coherente: cada refinamiento agrega información significativa sin contradecir el nivel anterior. En contraste, si la segmentación fuera arbitraria, la Purity no mostraría mejora consistente al refinar la granularidad.

\subsubsection{Comparación: ¿Cuántico vs. Clásico?}
\label{subsubsec:comparacion_quantico_clasico}

\begin{table}[Comparación de segmentación: HRP Clásico vs QHRP]{TB:COMP_SEGMENTATION}{Resumen de diferencias en estrategia de segmentación.}
\begin{tabular}{lll}
\hline
\textbf{Aspecto} & \textbf{HRP Clásico} & \textbf{QHRP Cuántico} \\
\hline
Base matemática & Correlación lineal & Matrices de densidad cuántica \\
\hline
Criterio de agrupamiento & Similitud en dendograma & Similitud en firmas cuánticas \\
\hline
Grupo más grande & 147 activos, Purity=0.129 & 31 activos, Purity=0.290 \\
\hline
Interpretación & Mezcla heterogénea & Familias de factores de riesgo \\
\hline
Robustez ante shocks & Baja (correlaciones colapsan) & Alta (factores de riesgo coherentes) \\
\hline
Rendimiento financiero & Sharpe agregado: 1.238 & Sharpe agregado: 1.444 \\
\hline
\end{tabular}
\end{table}

Como muestra la Tabla~\ref{TB:COMP_SEGMENTATION}, el resultado principal es que QHRP no solo obtiene mejor rendimiento financiero (Sharpe 1.444 vs. 1.238), sino que además produce segmentaciones más económicamente coherentes. Los grupos cuánticos responden a factores de riesgo interpretables (rentas variables, diversificadores, coberturas), mientras que los clásicos mezclan indiscriminadamente activos con dinámicas opuestas.

Esta coherencia sugiere que la representación cuántica está capturando estructura no lineal real de las dependencias entre activos, más allá de meras correlaciones de segundo orden.

\subsubsection{Implicaciones para optimización y diversificación}
\label{subsubsec:implicaciones_optimizacion_diversificacion}

Una segmentación coherente tiene implicaciones prácticas significativas:

\begin{enumerate}
	\item \textbf{Asignación de riesgo más estable:} si cada grupo representa un factor de riesgo diferente, la ponderación dentro del grupo es más robusta a cambios de régimen que si el grupo es heterogéneo.
	
	\item \textbf{Identificación de coberturas reales:} el aislamiento de materias primas y renta fija como grupo aparte permite construir subcarteras con objetivos específicos (ej., cartera de riesgo con renta variable, cartera de coberturas con diversificadores).
	
	\item \textbf{Reducción de error de estimación:} estimar covarianzas dentro de un grupo de 5-6 activos coherentes produce matrices más estables que estimar dentro de un mega-grupo de 147 activos heterogéneos.
\end{enumerate}

En conclusión, el análisis de segmentación sugiere que la ventaja de QHRP frente a HRP clásico no solo es empírica (mejor Sharpe), sino también estructural (mejor coherencia económica). Esta coherencia constituye una razón adicional para preferir el enfoque cuántico en contextos donde la estabilidad de la segmentación es crítica.

\subsection{Síntesis}
Los resultados muestran que no existe un método universalmente dominante: 
la elección depende del equilibrio deseado entre expresividad del modelo y coste computacional. 
El enfoque cuántico aporta una representación más rica de dependencias no lineales; 
el clásico aporta eficiencia extrema; y el kernel se sitúa como compromiso intermedio. 

Adicionalmente, el análisis de segmentación revela que QHRP produce agrupamientos coherentes con factores económicos reales (diversificadores, renta variable pura, coberturas), mientras que HRP clásico tiende a crear mega-grupos heterogéneos. Esta coherencia estructural, combinada con el mejor rendimiento financiero (Sharpe agregado 1.444 vs. 1.238), sugiere que la representación cuántica captura estructura significativa que trasciende la correlación lineal.

En consecuencia, la elección del método no es puramente técnica, sino dependiente del contexto operativo en el que se pretenda aplicar y del peso que se le dé a la interpretabilidad y robustez de la segmentación.
Esta lectura conjunta (rendimiento + coste + coherencia estructural) es la base para la discusión final del capítulo de conclusiones.

